\documentclass[11pt]{article}

\usepackage[utf8]{inputenc}
\usepackage{geometry}
\usepackage{graphicx}
\usepackage{xcolor}
\usepackage{setspace}
\usepackage{listings}
\usepackage{amsmath}
\usepackage{booktabs}
\usepackage{amssymb}
\usepackage{longtable}
\usepackage{array}
\usepackage{enumitem}
\usepackage{float}
\usepackage{hyperref}
\hypersetup{
    colorlinks=true,
    linkcolor=blue,
    urlcolor=blue,
    citecolor=blue
}

\setlist[itemize]{leftmargin=1.3em, noitemsep, topsep=0pt}
\setlist[enumerate]{leftmargin=1.3em, noitemsep, topsep=0pt}

\geometry{margin=1in}

\newcommand\sectionspath{}

\begin{document}

\vspace*{3cm}
\noindent
{\Huge \textcolor{blue}{\textbf{Tensor Omics (TOX)}}}\\[0.5cm]
\noindent
{\Huge \textcolor{blue}{\textbf{Analyze gene expression data and evolutionary relationships between genes}}}\\[0.5cm]
\vspace*{1cm}

\noindent
{\Huge \textcolor{blue}{\textbf{User’s Guide}}}

\vspace{2cm}

\noindent
\textbf{\large
Aaron Schröder,
Alexander Schwarzpaul,
Anna Beketova,
Asis Hallab,
Franz-Eric Sill,
Jitu Daba,
Jonathan Kurz,
Laszlo Lang,
Luka Faensen,
Mohamed Akdi,
Stephanie Ped,
Vivian Bass Vega
}\\[0.2cm]

{\small University of Applied Sciences Bingen, Germany }

\vspace{2cm}

\noindent
\textbf{Last revised \today}

\tableofcontents
\newpage

\input{tox_manual/methods_and_theory}
\section{Error Handling}
\label{sec:error-handling}

In case you encounter errors while using TOX, the following error codes may
help you identify the issue. The error codes are grouped into categories based
on their nature, such as I/O errors, format validation errors, memory errors, and
internal errors.

\begin{longtable}
    {>{\ttfamily}p{3cm} p{10cm}}
    \caption{TOX Error Codes Reference}
    \\ \hline \textbf{Error Code} & \textbf{Description} \\ \hline \endfirsthead

    \multicolumn{2}{c}%
    {{\bfseries \tablename\ \thetable{} -- continued from previous page}} \\ \hline
    \textbf{Error Code} & \textbf{Description} \\ \hline \endhead

    \hline \multicolumn{2}{r}{{Continued on next page}} \\ \endfoot

    \hline \endlastfoot

    0 & Success \\[6pt]

    \multicolumn{2}{l}{\textbf{1xx: I/O / File Reading Errors}} \\[3pt] 
    101 & Could not open file. \\ 
    102 & Could not read magic number. \\ 
    103 & Could not read array type code. \\ 
    104 & Could not read number of dimensions. \\ 
    105 & Could not read array dimensions. \\ 
    106 & Could not read character length. \\ 
    107 & Could not read array data. \\ 
    112 & Could not write magic number \\ 
    113 & Could not write array type code\\ 
    114 & Could not write number of dimensions \\
    115 & Could not write dimensions \\ 
    116 & Could not write character length \\
    117 & Could not write array data \\[6pt]

    \multicolumn{2}{l}{\textbf{2xx: Format / Input Validation Errors}} \\[3pt]
    200 & Invalid format detected (e.g., magic mismatch or corrupt format). \\ 
    201 & Invalid input provided. \\ 
    202 & Empty input arrays provided (e.g., zero-length arrays). \\ 
    203 & Dimension mismatch detected (shape inconsistencies). \\
    204 & NaN or Inf found in input data where not allowed. \\ 
    205 & Unsupported data type encountered (e.g., complex types). \\ 
    206 & Array size mismatch \\
    208 & Array index out of bounds \\ 
    209 & Division by zero detected \\[6pt]

    \multicolumn{2}{l}{\textbf{3xx: Memory Errors}} \\[3pt] 
    301 & Memory allocation failed. \\ 
    302 & Null pointer reference encountered. \\[6pt]

    \multicolumn{2}{l}{\textbf{5xxx: Fortran Runtime / Unit State Errors}} \\[3pt] 
    5002 & Fortran runtime error: unit not open or not connected. \\[6pt]

    \multicolumn{2}{l}{\textbf{9xxx: Internal / Unknown Errors}} \\[3pt] 9001 & Internal
    error: unexpected state or invariant violated. \\ 9999 & Unknown error (fallback
    for unexpected errors). \\
\end{longtable}

\section{Implementation Notes}

This section details the practical pipeline for using the TOX library to read and
process tabular data. The workflow is consistent across Fortran, Python, and R,
starting with reading a CSV file into a staging area of strings and then converting
columns to their proper data types.

\renewcommand{\sectionspath}[1]{tox_manual/fortran_pipeline/#1}

\lstset{ language=Fortran, basicstyle=\ttfamily\small, keywordstyle=\color{blue}, commentstyle=\color{green!40!black}, stringstyle=\color{purple}, showstringspaces=false, breaklines=true, frame=single, rulecolor=\color{black!30}, backgroundcolor=\color{black!5} }

\subsection{Fortran Pipeline}
\subsubsection{Overview}
\begin{enumerate}
    \item Data dependencies and relationships All arrays read have dependencies to
          each other. The gene ids array stores the ids, column number seven of the expression
          vectors belongs to entry number seven in the gene ids. Entry number nine
          in the gene to family mapping array, belongs to entry nine in the gene ids.
          These dependencies must be given at any time, therefore, only the provided
          functions shall be used to modify arrays in any form. For error code details
          please check \hyperref[sec:error-handling]{error handling section}.
          
    \item Error handling reference (placeholder for now)
          
    \item General usage principle (Placeholder for now)
\end{enumerate}

\input{\sectionspath{data_management}}
\input{\sectionspath{data_persistence}}
\input{\sectionspath{data_processing}}
\input{\sectionspath{family_based_analysis}}
\input{\sectionspath{rap_projection_and_analysis_functions}}
\input{\sectionspath{paralog_analysis}}
\input{\sectionspath{trajectory_contribution_analysis}}
\input{\sectionspath{fortran_utils}}
\renewcommand{\sectionspath}[1]{tox_manual/python_pipeline/#1}

\lstset{ language=Python }

\subsection{Python Pipeline}
\subsubsection{Overview}
The Python workflow is exposed through the \texttt{tensoromics\_functions.py}
module. The functions provide a high-level interface that mirrors the Fortran pipeline.

\input{\sectionspath{data_management}}
\input{\sectionspath{data_persistence}}
\input{\sectionspath{data_processing}}
\input{\sectionspath{family_based_analysis}}
\input{\sectionspath{rap_projection_and_analysis_functions}}
\input{\sectionspath{paralog_analysis}}
\input{\sectionspath{trajectory_contribution_analysis}}
\input{\sectionspath{python_utils}}
\input{tox_manual/r_pipeline}

\end{document}